\documentclass[a4paper,11pt]{ctexart}
\usepackage{xcolor}
\usepackage{array}
\usepackage{booktabs}
\usepackage{hyperref}
\hypersetup{colorlinks=true,linkcolor=blue!70!black}
\linespread{1.35}

\newcommand{\draft}{\textcolor{orange!80!black}{\textbf{【待写】}}}

\title{\textcolor{blue!70!black}{\Huge\textbf{临床医学模块 · 课程大纲}}}
\author{}
\date{}

\begin{document}
\maketitle
\thispagestyle{empty}

\section*{第零层 · 入门故事}

\begin{enumerate}
\item[\textbf{第1课}] \textbf{为什么打针比吃药见效快？——药怎么进到身体里} \draft \\
\textbf{内容：} 口服→消化→吸收→血液→目标（慢但方便）。注射→直接血液→目标（快但有创）。不同的给药途径：口服、注射、吸入、贴片、滴眼、栓剂。
\item[\textbf{第2课}] \textbf{医生怎么"看到"骨头里面的？——X射线的发现} \draft \\
\textbf{内容：} 伦琴发现X射线—CT（多张X光拼成立体）、B超（用声音不用辐射）、MRI（用磁场不用射线）。情景：为什么骨折了先拍X光而不是CT？
\item[\textbf{第3课}] \textbf{疫苗是怎么工作的？——让免疫系统提前"认识"敌人} \draft \\
\textbf{内容：} 灭活疫苗/减毒疫苗/mRNA疫苗/腺病毒载体疫苗。mRNA疫苗为什么能"设计"出来只用了一年？（之前的技术积累用了20年）。情景：新冠疫苗的快速研发是"科学奇迹"还是"早有准备"？
\end{enumerate}

\section*{深入课程}

\begin{enumerate}
\item[\textbf{第1课}] \textbf{人体解剖与生理——身体的地图} \draft \\
\textbf{内容：} 八大系统。标志性解剖结构（心/肺/肝/肾/脑的位置和功能）。体腔、骨性标志、重要血管神经。情景：为什么阑尾炎是右下腹痛？
\item[\textbf{第2课}] \textbf{病理生理——疾病是怎么发生的？} \draft \\
\textbf{内容：} 炎症（红/肿/热/痛/功能障碍的分子机制）、肿瘤（良性vs恶性、转移的分子途径）、缺血再灌注损伤。情景：为什么心梗要"争分夺秒"？
\item[\textbf{第3课}] \textbf{药理学——药为什么有用也有副作用？} \draft \\
\textbf{内容：} 受体-药物相互作用、药代动力学（ADME：吸收/分布/代谢/排泄）、剂量-反应曲线、治疗指数。情景：为什么对乙酰氨基酚（扑热息痛）在安全剂量内很安全，但超量就会肝衰竭？
\item[\textbf{第4课}] \textbf{诊断技术——从听诊器到基因测序} \draft \\
\textbf{内容：} 物理诊断（视触叩听）、影像诊断（X光/CT/MRI/超声/PET）、实验室诊断（血/尿/生化/免疫）。基本原理、各自优劣势。情景：为什么CT是"断面图"而MRI可以显示软组织细节？
\item[\textbf{第5课}] \textbf{外科手术——从切开到缝合} \draft \\
\textbf{内容：} 无菌原则（感染是外科的头号敌人）、麻醉（全麻/局麻/椎管麻醉）、手术入路的选择。微创手术（腹腔镜/机器人手术）vs 开放手术。情景：达芬奇手术机器人的"力反馈"为什么做不出来？
\item[\textbf{第6课}] \textbf{北京在地——北京的顶级医院与医学研究} \draft \\
\textbf{内容：} 协和医院（全国疑难重症中心）、北医三院（运动医学）、北京医院（老年医学）、天坛医院（神经外科）、阜外医院（心血管）、肿瘤医院。海淀区的"医学谷"（北大医学部+中科院+生物医药产业）。
\end{enumerate}

\end{document}
